% Source: https://artofproblemsolving.com/community/c3h2168951_2020_mock_amc_11_problem_7_11_discussion
% Stable identity: aops-topic-2168951-q07
Put Carol at $C=(0,0)$, the room at $0\le x,y\le3$ and the pillar at $1\le x,y\le2$. A point is hidden when the segment from $C$ to it passes through the pillar. The rays from $C$ that touch the pillar lie between the rays through $(2,1)$ and $(1,2)$, i.e. between the lines $y=\dfrac {x}{2}$ and $y=2x$; these rays meet the walls at $\left(3,\dfrac{3}{2}\right)$ and $\left(\dfrac{3}{2},3\right)$. The hidden region is the part of this wedge beyond the far sides of the pillar (shaded light grey).

\begin{asy}
usepackage("amsmath");
size(5cm);
defaultpen(fontsize(10pt));
fill((2,1)--(3,1.5)--(3,3)--(1.5,3)--(1,2)--(2,2)--cycle,lightgray);
fill((1,1)--(2,1)--(2,2)--(1,2)--cycle,gray);
draw((0,0)--(3,0)--(3,3)--(0,3)--cycle,linewidth(1));
draw((1,1)--(2,1)--(2,2)--(1,2)--cycle,linewidth(0.9));
draw((0,0)--(3,1.5),red+dashed); draw((0,0)--(1.5,3),red+dashed);
dot((0,0)); label("$C$",(0,0),SW);
label("$(2,1)$",(2,1),SE,fontsize(8pt)); label("$(1,2)$",(1,2),W,fontsize(8pt));
label("$(3,\dfrac{3}{2})$",(3,1.5),E,fontsize(8pt)); label("$(\dfrac{3}{2},3)$",(1.5,3),N,fontsize(8pt));
label("pillar",(1.5,1.5),white+fontsize(8pt));
label("hidden",(2.45,2.45),fontsize(8pt));
label("\mbox{}",(min(currentpicture,user=true)+max(currentpicture,user=true))/2,black); // keeps the text after the figure black when LaTeX typesets the labels inline
\end{asy}

The part of the wedge inside the room is the quadrilateral $C\left(3,\dfrac{3}{2}\right)(3,3)\left(\dfrac{3}{2},3\right)$; the diagonal from $C$ to $(3,3)$ splits it into two triangles of area $\dfrac{1}{2}\cdot\dfrac{3}{2}\cdot3=\dfrac{9}{4}$ each, so its area is $\dfrac{9}{2}$. The part of the wedge in front of the far sides of the pillar is the quadrilateral $C(2,1)(2,2)(1,2)$, made of two triangles of area $\dfrac{1}{2}\cdot1\cdot2=1$, so its area is $2$ (it includes the pillar). The hidden floor has area $\dfrac{9}{2}-2=\dfrac{5}{2}$.

The floor has area $9-1=8$, of which Carol sees $8-\dfrac{5}{2}=\dfrac{11}{2}$. The fraction she sees is
\[\dfrac{11}{2}\div8=\boxed{\dfrac{11}{16}},\]
answer $\textbf{(C)}$.
